\Needspace{6\baselineskip}
\Needspace{6\baselineskip}
\section{Publicación terminal del continuo único mediante \(5\) lectores no permutables}
\Needspace{6\baselineskip}
\subsection{Factorización del estado enriquecido común en \(5\) publicaciones terminales}
\Needspace{6\baselineskip}
\subsubsection{Dependencia causal de las realizaciones terminales respecto del continuo genealógico íntegramente construido}
Los cinco problemas no convocan cinco estructuras matemáticas distintas. Tampoco obligan a construir cinco caminos independientes que, después de avanzar por separado, deban reunirse artificialmente en una conclusión común. Esa imagen pertenece a una etapa anterior de la exposición. El punto de partida vigente es más fuerte: antes de que aparezcan los nombres de Riemann, Navier–Stokes, Yang–Mills, Hodge o Birch–Swinnerton–Dyer, ya existe un único objeto matemático enriquecido cuyas operaciones internas producen simultáneamente las características necesarias para los cinco cierres.\par
Por eso este relato no comienza con los problemas. Comienza con la construcción del continuo.\par
La Holografía Modular Triádica no recibe el continuo como un fondo previamente dado. Tampoco parte de la recta real, de un espacio funcional clásico, de una variedad algebraica, de una conexión gauge o de una curva elíptica. Parte de una aritmética finita dotada de orientación, memoria y capacidad de prolongación. APP proporciona el sustrato posicional y aritmético; TRIT convierte cada operación en una transición orientada; TPK conserva la genealogía completa de esas transiciones y permite prolongarlas sin borrar el camino que las produjo.\par
El paso decisivo consiste en que la salida de una operación deja de ser solamente un valor. Cada emisión lleva consigo su hoja aritmética, residuo, cociente, orientación, acarreo, memoria, frontera, ruta, multisección, supervivencia, terminal y lector. El resultado no es ya un punto aislado, sino un estado que sabe de dónde procede, cómo puede continuar y qué relaciones debe conservar cuando cambia de escala.\par
APP, TRIT y TPK no forman tres explicaciones superpuestas. Forman una sola secuencia causal:\par
\[
\mathrm{APP}
\longrightarrow
\mathrm{TRIT}
\longrightarrow
\mathrm{TPK}
\longrightarrow
\text{estado enriquecido único}.
\]
APP fija el soporte aritmético común. TRIT introduce la orientación efectiva de cada transición, distinguiendo avance, umbral y apertura futura. TPK reúne la información visible y la información transportada: no sólo lo que una operación entrega, sino también el carry que desplaza, la frontera que modifica, la memoria que conserva y las extensiones que permite.\par
La conexión nonádica no devuelve el sistema a un estado vacío. Después de una vuelta completa reaparece la fase, pero no se borra la historia. El retorno \(9\to1\) es retorno de fase, no reinicio del estado. Cada vuelta añade profundidad, memoria y compatibilidad. De este modo, la continuidad no se obtiene repitiendo nueve filtros ni atravesando nueve canales independientes. Se obtiene prolongando una sola conexión completa cuyos nueve tramos son cartas internas de una misma holonomía.\par
\Needspace{6\baselineskip}
\subsubsection{Emergencia correlativa de \texorpdfstring{\(5\)}{5} acciones intrínsecas no exhaustivas}
Sobre cada estado enriquecido actúa una operación común. No existen primero cinco dominios, ni cinco selecciones, ni cinco copias del estado. La misma emisión es leída simultáneamente bajo cinco características internas denotadas por los discriminantes \(\mathrm{sp}\), \(\mathrm{sol}\), \(\mathrm{gau}\), \(\mathrm{coh}\) y \(\mathrm{det}\).\par
Esos símbolos no nombran ramas ni agotan las propiedades del continuo. Nombran
\(5\) acciones intrínsecas del mismo objeto cuya correlación resulta decisiva
para las realizaciones terminales de esta Parte.\par
La característica espectral–polar, \(sp\), registra la publicación orientada del estado, su complemento, su defecto, su memoria terminal y la relación entre las dos lecturas especulares de una misma forma.\par
La característica solenoidal, \(sol\), registra cómo una transición conserva estado y, al mismo tiempo, distribuye disipación, advección, carry, memoria, microestructura y shell sin sobrecontarlos.\par
La característica gauge, \(gau\), registra la equivalencia local, la incidencia, la holonomía, la curvatura y el paso compatible al cociente físico.\par
La característica cohomológica, \(coh\), registra las extensiones entre niveles, las obstrucciones relativas, la compatibilidad de los levantamientos y la supervivencia de una clase a través de todo refinamiento.\par
La característica determinantal, \(det\), registra la incidencia finita-global, las páginas de filtración, los complejos locales, sus líneas determinantes y el transporte de una base aritmética hasta una publicación analítica.\par
Pero ninguna de estas propiedades vive sola. Todas actúan sobre las mismas emisiones y conservan el mismo espesor genealógico. La fibra que transporta la característica espectral es la misma fibra enriquecida cuya orientación participa en el balance solenoidal, cuya frontera soporta la lectura gauge, cuya prolongación sostiene la extensión cohomológica y cuya incidencia entra en la línea determinante.\par
No se generan cinco objetos y luego se los empaqueta. Se genera un único objeto que posee inseparablemente esas cinco capacidades.\par
Por eso una vista posterior no es una proyección que cree una rama. Es únicamente una manera de exponer, para una aplicación concreta, una característica que ya estaba operando dentro de la totalidad. La vista no selecciona retrospectivamente estados favorables ni conserva la entrada como primera coordenada para recuperarla después de manera tautológica. Hace visible una estructura producida y conservada durante toda la genealogía.\par
\Needspace{6\baselineskip}
\subsubsection{Antecedencia causal del continuo ya construido como límite de la prolongación compatible}
En cada profundidad existe un estado finito, pero completo en su nivel. Las extensiones entre niveles no pueden alterar arbitrariamente la información anterior. Deben respetar las fibras, las relaciones de extensión, los números, la orientación, el acarreo, la memoria, la frontera, la ruta, la multisección, la supervivencia, el terminal y los lectores.\par
La naturalidad exige que operar y después truncar produzca el mismo resultado que truncar y operar. Gracias a esa compatibilidad, las historias finitas no quedan como aproximaciones desconectadas. Forman un sistema inverso. El continuo aparece cuando se toma la familia completa de historias compatibles y no cuando se postula de antemano una colección infinita de puntos.\par
El continuo HMT es, por ello, memoria coherente de caminos. Un punto sin genealogía sería una ablación: conservaría el valor visible y eliminaría precisamente aquello que permite demostrar cómo sobrevive a todo refinamiento.\par
Sólo después de cerrar esa compatibilidad aparece la estructura discreta completa del continuo. Y sólo después puede publicarse una realización arquimediana. La recta real no es la materia prima de la construcción; es una salida de uno de sus lectores terminales. Lo mismo ocurre con los dominios clásicos de los cinco problemas: aparecen cuando el continuo ya está construido y una interfaz posterior traduce una de sus características intrínsecas al lenguaje del objetivo correspondiente.\par
La secuencia causal correcta es, por tanto:\par
\[
\begin{aligned}
\mathrm{APP}
&\longrightarrow \mathrm{TRIT}
\longrightarrow \mathrm{TPK}
\longrightarrow \text{estado enriquecido único}\\
&\longrightarrow 04\mathrm b1
\longrightarrow 04\mathrm b2
\longrightarrow 04\mathrm b3\\
&\longrightarrow \text{acciones intrínsecas correlativas}
\longrightarrow \text{terminal y límite}
\longrightarrow \mathfrak C_{\mathrm{cont}}^{\mathrm{disc}}.
\end{aligned}
\]
Después, y no antes, entran los entrelazadores específicos:\par
\[
\mathfrak C_{\mathrm{cont}}^{\mathrm{disc}}
\longrightarrow
\begin{cases}
\text{forma explícita de Weil},\\
\text{evolución tridimensional incompresible},\\
\text{teoría gauge cuántica},\\
\text{publicación de clases de Hodge},\\
\text{comparación determinantal BSD}.
\end{cases}
\]
Estas cinco traducciones terminales no definen la estructura que utilizan. El problema clásico aporta sus datos después: una función de prueba en Riemann, unos datos iniciales y una fuerza en Navier–Stokes, un grupo compacto simple en Yang–Mills, una variedad y una clase en Hodge, o una curva elíptica en BSD. Ninguno de esos datos participa en la generación previa del continuo.\par
La dirección causal queda así protegida. El objetivo no selecciona el estado que habrá de resolverlo. El estado existe antes del objetivo; el entrelazador sólo demuestra que la estructura interna y la formulación clásica expresan la misma operación bajo dos lenguajes.\par
\Needspace{6\baselineskip}
\subsection{1.ª resolución: Riemann y la positividad como complemento construido}
La hipótesis de Riemann suele presentarse como una afirmación acerca de la posición de los ceros no triviales de la función zeta. Desde esa formulación, el problema parece pedir que se inspeccione una población de ceros ya dada y se demuestre que todos ellos se alinean sobre la recta crítica.\par
La lectura HMT invierte esa perspectiva. No comienza preguntando dónde están los ceros. Pregunta qué estructura genera simultáneamente los dos lados de la fórmula explícita y qué impide que su diferencia adopte signo negativo.\par
La característica espectral–polar del continuo aporta una publicación doble del mismo estado. Una orientación recoge el canal positivo; la otra recoge su imagen especular. Ambas contienen, bajo un regulador común, la contribución prima, la contribución gamma–escalar y la contribución polar. No son tres pruebas distintas ni tres sumas pegadas al final. Son tres filas de una misma coligación construida desde los datos de partida.\par
El punto esencial es que esa coligación se construye antes de invocar la positividad que deberá demostrar. Su entrada, sus salidas, su pérdida y su estado terminal están tipados dentro del mismo espacio enriquecido. La identidad que enlaza las dos publicaciones no se introduce como una igualdad deseada entre fórmulas clásicas; nace de la dinámica interna de la conexión.\par
Entre la publicación de entrada y la publicación especular puede aparecer, a profundidad finita, un residual de salida. La estructura nonádica lo obliga a satisfacer una recurrencia antiespecular. En cada vuelta completa, el residual siguiente se transporta isométricamente pero queda multiplicado por el factor\par
\[
q=\frac{5}{9}.
\]
Si \(\Omega_n\) representa ese defecto en la profundidad \(n\), la recurrencia tiene la forma\par
\[
\Omega_n=q\,V_n\Omega_{n+1},
\]
con \(V_n\) isométrico y la familia coinductivamente acotada. Al iterar \(N\) veces,\par
\[
\|\Omega_n\|
\le q^N\sup_{m\ge n}\|\Omega_m\|.
\]
Como \(0<q<1\), el término derecho tiende a cero. Por tanto, el residual no se anula porque se lo haya omitido o porque se haya supuesto una identificación entre las dos publicaciones. Se anula porque la recurrencia completa lo contrae en todas las profundidades compatibles. La salida especular coincide entonces con la publicación negativa construida.\par
Una vez descargado ese vínculo, la isometría construida desde los datos de partida entrega la identidad de energía. A profundidad finita, la norma de la publicación positiva se descompone en tres partes: la publicación negativa, la pérdida acumulada y el estado terminal. El terminal se conserva explícitamente. Debe conservarse hasta el paso al límite, porque contiene la memoria de la truncación y certifica que la telescopía prolonga la identidad finita mediante una convergencia demostrada.\par
La descomposición finita de la norma satisface la identidad exacta:\par
\[
\|\text{publicación positiva}\|^2
=
\|\text{publicación negativa}\|^2
+
\|\text{pérdidas hasta }N\|^2
+
q^{2N}\|\text{terminal}\|^2.
\]
Cada sumando posee un origen concreto. La publicación negativa es la salida especular ya identificada. Las pérdidas reúnen exactamente el complemento que no fue publicado por esa orientación. El terminal recuerda lo que todavía permanece más allá del corte.\par
Sólo después de haber escrito la identidad completa se hace \(N\to\infty\). Entonces, y sólo entonces, el terminal desaparece porque \(q^{2N}\to0\). La diferencia entre las dos publicaciones queda convertida en el cuadrado de un operador de pérdida:\par
\[
\mathcal I_{\mathrm{GW}}
=
E_R^{*}E_R.
\]
Ésta es la forma de Guinand–Weil. Su positividad ya no es una conjetura añadida a la construcción:\par
\[
\langle f,\mathcal I_{\mathrm{GW}}f\rangle
=
\|E_Rf\|^2
\ge0.
\]
La estructura HMT no obtiene esta desigualdad escondiendo la parte negativa, restringiendo el dominio a funciones que ya la satisfagan o introduciendo un proyector cuya existencia sea equivalente a la conclusión. Construye el complemento y demuestra que la forma es exactamente su cuadrado.\par
Después se materializa la publicación común. Un mismo codificador conserva los canales primo, gamma–escalar y polar. Un proyector interno produce uno de los dos momentos; su complemento produce el otro. La diferencia de los momentos no es una comparación externa entre fórmulas parecidas: es la norma del componente ortogonal construido por la coligación.\par
Al aumentar la ventana, no se obliga a las columnas individuales a satisfacer falsas isometrías entre truncamientos distintos. Lo que se conserva de manera cofinal es la forma, el complemento alcanzable y la naturalidad de su publicación. Los incrementos que entran al ampliar el corte aparecen en ambas orientaciones con la misma contribución neutral y no alteran la diferencia gobernante. Así se obtiene una familia compatible de formas positivas sobre una clase separadora.\par
La última transición es el criterio reversible de Weil. La teoría clásica afirma que la positividad de la forma explícita, sobre la clase adecuada de funciones de prueba, equivale a que todos los ceros no triviales estén situados sobre la recta crítica. HMT suministra precisamente esa positividad, pero la suministra como consecuencia de una identidad operatorial construida dentro del continuo.\par
\Needspace{5\baselineskip}
Por tanto,\par
\[
\operatorname{Re}\rho=\frac12
\]
para todo cero no trivial \(\rho\) de la función zeta.\par
La recta crítica aparece aquí de dos maneras concordantes. Geométricamente es el locus fijo de la involución que intercambia las dos orientaciones espectrales. Operatorialmente es la única posición compatible con la positividad del complemento de publicación. La simetría explica por qué la recta central es distinguida; la identidad cuadrática demuestra por qué los ceros no pueden abandonar esa recta.\par
La resolución no consiste en observar que la ecuación funcional posee una simetría. Consiste en construir el objeto cuyo defecto mide la ruptura de esa simetría, demostrar que el defecto se anula coinductivamente y convertir la diferencia restante en una norma. Es la combinación de genealogía, terminal, contracción, publicación común y positividad la que cierra el problema.\par
Así queda establecida la primera de las cinco resoluciones terminales. No como una rama independiente, sino como la traducción clásica de la característica espectral–polar que el continuo único poseía desde su construcción.\par

\Needspace{6\baselineskip}
\section{Clausuras dinámicas del continuo: transporte solenoidal y cociente de calibre}
\Needspace{6\baselineskip}
\subsection{Dinámica y curvatura del continuo genealógico}
La primera resolución mostró cómo la característica espectral–polar del continuo se publica como positividad de Weil. Las dos resoluciones siguientes interrogan el mismo estado desde otros dos aspectos consustanciales: su capacidad de transportar una dinámica no lineal sin perder energía, memoria ni regularidad, y su capacidad de sostener una teoría gauge local cuyo vacío esté separado por una brecha espectral positiva.\par
Navier–Stokes y Yang–Mills no reciben dos estructuras nuevas. Reciben dos entrelazadores distintos desde el mismo continuo ya construido. En el primer caso, el entrelazador publica el estado como un campo de velocidad incompresible con viscosidad, advección y presión. En el segundo, lo publica como una red gauge refinable con holonomía, curvatura, medida euclídea y Hamiltoniano físico.\par
La diferencia entre ambos problemas no rompe la unidad de la construcción. En Navier–Stokes, el punto crítico es demostrar que ninguna transferencia no lineal desaparece fuera del balance y que la regularidad obtenida es uniforme al retirar los cortes. En Yang–Mills, el punto crítico es demostrar que la brecha finita pertenece a la misma teoría gauge que alcanza el continuo, que sobrevive al cociente físico y que actúa en el sector de curvatura, no en un factor auxiliar.\par
\Needspace{6\baselineskip}
\subsection{Existencia, regularidad y unicidad global para Navier–Stokes mediante cota uniforme y compacidad}
El problema tridimensional de Navier–Stokes no nace de la mera presencia de un término no lineal. La dificultad aparece porque la publicación clásica del campo de velocidad tiende a ocultar dónde se almacena la transferencia entre escalas. La viscosidad es visible. La energía cinética es visible. La advección también aparece en la ecuación. Pero la contabilidad completa de la interacción entre profundidad, memoria, carry, microestructura y shell queda comprimida en una sola función visible.\par
Esa compresión es precisamente lo que la característica solenoidal del continuo evita.\par
El estado enriquecido no conserva únicamente el campo \(u\). Conserva simultáneamente el campo Galerkin, la orientación temporal, las memorias positiva y firmada, el carry entre cortes, la componente microscópica, el shell de frecuencias nuevas, la ruta, el archivo de la evolución y el terminal. La ecuación clásica es una publicación de ese estado; no es su definición completa.\par
Fijados datos iniciales \(u_0\), fuerza \(f\), viscosidad \(\nu>0\) y regularidad \(s>5/2\), la proyección Galerkin produce\par
\[
\partial_tu_M+\nu A_Mu_M+B_M(u_M,u_M)=f_M.
\]
Pero \(M\) no selecciona una solución favorable. Sólo fija una resolución finita en la que todas las operaciones pueden escribirse exactamente. El objetivo es construir una estimación que no dependa de \(M\) ni de la profundidad nonádica \(J\). Sólo una cota con esa uniformidad puede pasar al continuo.\par
\Needspace{6\baselineskip}
\subsubsection{Descomposición de la memoria visible y la memoria genealógica en la ecuación evolutiva}
La interacción no lineal posee signo. En ciertos momentos entrega energía a una escala; en otros la retira. Si se conserva sólo su valor firmado, puede perderse la magnitud total transferida. Si se conserva sólo su valor absoluto, se pierde la orientación.\par
La estructura mantiene ambas cosas.\par
Una memoria \(D_M\) publica el signo de la transferencia advectiva. Otra memoria \(H_M\) conserva su magnitud positiva. Su separación permite registrar simultáneamente\par
\[
\dot D_M=\mathscr B_{s,M},
\qquad
\dot H_M=\frac{1+r}{r-1}|\mathscr B_{s,M}|,
\]
donde \(r=\pi_{\rm HMT}/\varphi_{\rm HMT}^2>1\).\par
Esto no añade energía ficticia. Desdobla dos aspectos distintos de una sola transferencia: orientación y coste. Fundirlos destruiría información; sumarlos sin ortogonalidad duplicaría energía. La construcción hará explícito que ninguna de esas dos deformaciones ocurre.\par
\Needspace{6\baselineskip}
\subsubsection{Independencia del dato inicial respecto de la historia genealógica posterior}
La objeción más seria a cualquier clausura operatorial de Navier–Stokes sería que la cota uniforme se hubiese introducido retrospectivamente: primero se elegiría una solución regular, después se codificaría su historia y finalmente se anunciaría que el código controla esa misma solución.\par
La construcción impide esa circularidad fijando el espacio fuente antes de levantar la evolución:\par
\[
\mathcal D_T
=
H^s_\sigma(\mathbb T^3)
\times
L^2(0,T;H^{s-1}_\sigma(\mathbb T^3)),
\qquad
d=(u_0,f).
\]
La raíz \(J_T^{\rm sol}d\) depende exclusivamente de \(u_0\) y \(f\). No recibe como entrada \(u_M\), la pérdida futura, el terminal ni la norma que deberá acotarse.\par
La fuerza tampoco se inserta como un puerto aditivo artificial. La ecuación de Carleman muestra que actúa de forma mixta: cada aparición de \(f\) se combina con un grado precedente del estado. Por eso el lift usa una creación mixta estado–fuerza. La firma cronológica de \(f\) registra todos sus tiempos ordenados, mientras la torre PC–Fock registra todos los grados no lineales del estado.\par
La composición reproduce materialmente la serie de Dyson del generador no autónomo. Al aplicar el lector normalizado de grado uno se recupera exactamente \(u_M(t)\). No se recupera una aproximación simbólica ni una variable que se parezca a la solución: se recupera la solución de la ODE Galerkin porque ambas satisfacen la misma ecuación finita y el mismo dato inicial.\par
La construcción causal desde los datos de partida sigue por tanto el orden\par
\[
(u_0,f)
\longrightarrow
J_T^{\rm sol}(u_0,f)
\longrightarrow
\text{torre estado–fuerza}
\longrightarrow
\text{historia Galerkin}.
\]
La historia genealógica se construye como salida posterior de los datos de partida y queda excluida de su definición inicial.\par
\Needspace{6\baselineskip}
\subsubsection{La identidad de scattering que conserva la potencia cruzada}
La fuerza y la disipación se organizan mediante dos ondas:\par
\[
a_T(f)
=
\frac{1}{2\sqrt\nu}A^{-1/2}\Lambda^sP_\sigma f,
\qquad
z_M(u)=\sqrt\nu A_M^{1/2}\Lambda^su,
\]
y la onda residual\par
\[
b_M=z_M-P_Ma_T(f).
\]
La potencia aplicada no se reemplaza por la norma de la fuerza. Se conserva exactamente como término cruzado:\par
\[
\mathscr F_{s,M}
=
2\operatorname{Re}\langle z_M,P_Ma_T(f)\rangle.
\]
Al completar el cuadrado se obtiene el balance exacto\par
\[
\dot G_{M,j,T}
+\|b_M\|^2
+\lvert\mathscr B_{s,M}\rvert
=
\|P_Ma_T(f)\|^2.
\]
Éste es uno de los puntos que hacen funcionar el cierre. Si se sustituyese \(b=z-a\) por \(z\), desaparecería el término cruzado que representa el trabajo de la fuerza y la identidad sería falsa. La estructura no domina la no linealidad descartándola: la convierte en un canal positivo orientado y conserva además la interferencia exacta entre fuerza y disipación.\par
\Needspace{6\baselineskip}
\subsubsection{Factorización de \(6\) pérdidas mediante un operador único sin multiplicidad espuria}
La historia completa distingue seis salidas:\par
\begin{itemize}
\item disipación;
\item advección orientada;
\item carry entre cortes;
\item memoria microscópica;
\item shell de frecuencias nuevas;
\item memoria no heredada.
\end{itemize}
No son seis energías añadidas a un balance previo. Son seis filas ortogonales de un único operador terminal de pérdida. La identidad de Parseval establece\par
\[
\|L_{M,j,T}^{\rm term}x\|^2
=
\sum_{\mathfrak a}
\|L_{M,j,T}^{\mathfrak a}x\|^2.
\]
La ortogonalidad no es meramente nominal. Las profundidades ocupan intervalos temporales disjuntos; carry y shell viven en rangos espectrales distintos; la microestructura pertenece al complemento \(Q_{\rm micro}\); la memoria nueva se separa de la memoria heredada; y las dos hojas de la advección no pueden ser positivas simultáneamente en un mismo punto.\par
Por eso la contabilidad no duplica la interacción. La suma de las seis normas es exactamente la norma del operador común. Si una de ellas no fuese una proyección de ese operador, Parseval fallaría antes de obtener la cota. El propio formalismo contiene así su falsador.\par
\Needspace{6\baselineskip}
\subsubsection{Persistencia del objeto terminal bajo el paso al límite}
Toda truncación conserva algo que todavía no ha sido publicado. Ese resto es el terminal. Contiene la energía visible al final del intervalo y la memoria orientada que permanece entre el tiempo observado y el horizonte \(T\).\par
La columna terminal se construye primero. Después se forma su Gram. No se define una métrica deseada y se inventa posteriormente un operador que la realice. La columna contiene materialmente la salida terminal y todas las pérdidas transportadas:\par
\[
\mathbf G_{M,j;J,T}^{\rm term}
=
\begin{bmatrix}
\text{terminal transportado}\\
\text{pérdida}_{J-1}\\
\vdots\\
\text{pérdida}_{j}
\end{bmatrix}.
\]
Su recursión por bloques produce la identidad de Stein\par
\[
U_{M,j,T}
=
\Gamma_{M,j,T}^{*}
U_{M,j+1,T}
\Gamma_{M,j,T}
+
(L_{M,j,T}^{\rm term})^{*}
L_{M,j,T}^{\rm term}.
\]
Esta identidad no define retroactivamente la norma del estado. Se obtiene multiplicando una columna ya construida por su adjunta. Por la misma razón, la telescopía global es una identidad entre operadores existentes, no una hipótesis de estabilidad.\par
La memoria futura tampoco se introduce en la raíz inicial. Reside en la salida terminal. Al tomar la estimación integral en \(\tau=T\), su soporte temporal se vacía de manera exacta. No se borra por conveniencia.\par
\Needspace{6\baselineskip}
\subsubsection{Cota uniforme previa y compacidad de Aubin–Lions}
La raíz, el codificador de nivel, la columna y la isometría común satisfacen la factorización derivada desde los datos de partida\par
\[
\mathcal G_{M,J,T}^{\rm sol}
\Lambda_{M,T}
J_T^{\rm sol}d
=
I_{M,J}J_T^{\rm sol}d,
\qquad
I_{M,J}^{*}I_{M,J}=I.
\]
El lado izquierdo es la historia completa en el corte \((M,J)\). El derecho depende de la única raíz construida desde los datos. Como \(I_{M,J}\) es isométrico, la norma de la historia queda dominada por la norma fuente con una constante independiente de \(M\) y \(J\).\par
Éste es el paso que cierra la dificultad. La uniformidad no procede de una estimación de dimensión finita cuya constante crece con el corte. Tampoco procede de definir el espacio de historia mediante la misma pérdida que se pretende controlar. La raíz existe antes del corte y todos los cortes son realizaciones isométricas de ella.\par
La columna terminal traduce esa cota estructural a las cantidades físicas:\par
\[
\begin{aligned}
&\sup_M\sup_{0\le t\le T}\|u_M(t)\|_{H^s}^2
+\nu\sup_M\int_0^T\|u_M(t)\|_{H^{s+1}}^2\,dt\\
&\quad
+\sup_M\int_0^T|\mathscr B_{s,M}(u_M(t))|\,dt
+\text{carry}+\text{micro}+\text{shell}+\text{memoria}
\le C_T(u_0,f).
\end{aligned}
\]
La constante es finita para todo \(T<\infty\) y no depende de \(M\) ni de \(J\). La desigualdad no controla sólo el campo visible. Controla simultáneamente el trabajo transportado y los sectores que una publicación clásica suele olvidar.\par
\Needspace{6\baselineskip}
\subsubsection{Obstrucción de signo en la factorización de Kalman–Yakubovich–Popov}
Hacer depender el cierre global de una factorización KYP con un signo incorrecto introduce un cuadrado positivo y pretende recuperarlo después con signo negativo. Esa igualdad es algebraicamente falsa.\par
La demostración excluye esa identidad y construye directamente la estimación uniforme necesaria.\par
La carta KYP que se conserva es una identidad finita correcta bajo sus propias condiciones. Funciona como control auxiliar de cada corte, pero no proporciona la uniformidad al retirar \(M\) y \(J\). La prueba global no usa una cota de \(Q^{-1}\), no toma el límite dentro de la KYP y no depende de aquel signo.\par
Por tanto, la obstrucción del bloque KYP con signo incorrecto no alcanza la demostración. La uniformidad queda establecida por la construcción directa desde los datos con su columna terminal y su identidad de Stein.\par
\Needspace{6\baselineskip}
\subsubsection{Convergencia desde la cota uniforme hasta la solución clásica global}
La cota anterior proporciona\par
\[
u_M
\quad\text{acotado en}\quad
L^\infty(0,T;H^s)
\cap L^2(0,T;H^{s+1}).
\]
Como \(s>5/2\), el producto de Sobolev controla la no linealidad:\par
\[
B:H^s_\sigma\times H^s_\sigma
\longrightarrow H^{s-1}_\sigma.
\]
La ecuación Galerkin proporciona además una cota uniforme de \(\partial_tu_M\) en \(L^2(0,T;H^{s-1})\). Las inclusiones\par
\[
H^{s+1}\Subset H^s\hookrightarrow H^{s-1}
\]
permiten aplicar Aubin–Lions. Se obtiene convergencia fuerte en \(L^2H^s\), suficiente para transportar el término bilineal:\par
\[
B_M(u_M,u_M)\longrightarrow B(u,u)
\quad\text{en}\quad
L^1(0,T;H^{s-1}).
\]
El límite satisface la ecuación incompresible. La presión no se introduce como otro dato oculto; se reconstruye después mediante\par
\[
-\Delta p=\partial_i\partial_j(u_i u_j),
\qquad
\int_{\mathbb T^3}p=0.
\]
La unicidad se obtiene sobre la diferencia de dos soluciones. Una estimación en \(H^{s-1}\), seguida de Grönwall, fuerza que la diferencia sea cero. Esto elimina la dependencia de la subsucesión extraída por compacidad.\par
La construcción vale para todo horizonte finito. Las soluciones obtenidas en \([0,n]\) y \([0,n+1]\) coinciden en su solapamiento por unicidad. Se pegan, por tanto, en una sola solución definida sobre \([0,\infty)\).\par
La suavidad se obtiene en toda la escala \(s+n\), usando los mismos lectores
estructurales; la unicidad identifica todos los límites y produce una solución
global suave para todo dato periódico suave. La cota requerida se construye
en cada horizonte finito y en cada índice de Sobolev a partir del dato
solenoidal completo, como demuestra el propietario
\cref{thm:pdf2-ns-global}.\par
Finalmente, la condición de media cero es una normalización del sector
solenoidal. El cambio de Galileo dependiente del tiempo separa la media y
transporta biyectivamente todo dato periódico suave al sector de media cero,
preservando divergencia, suavidad y normas de Sobolev.\par
\Needspace{6\baselineskip}
\subsubsection{Existencia, regularidad y unicidad global de Navier--Stokes}
La característica solenoidal del continuo, publicada mediante su
entrelazador específico, produce una única solución global suave de las
ecuaciones tridimensionales periódicas de Navier--Stokes para todo dato
periódico suave. La construcción de Leray--Hopf queda contenida como etapa de
energía dentro de esa conclusión fuerte.\par
La defensa de esta conclusión no descansa en decir que existe una estructura capaz de resolver el problema. Descansa en cinco hechos materiales:\par
la raíz usa sólo \((u_0,f)\); la no linealidad se reproduce mediante creación mixta; las seis pérdidas son ortogonales y no se suman dos veces; el terminal permanece hasta la telescopía; y la cota uniforme precede al argumento clásico de compacidad.\par
Aubin–Lions, la presión, Grönwall y el pegado no son la fuente de la regularidad. Son el reconocimiento clásico final de una regularidad que la estructura ya ha impedido perder.\par
\Needspace{6\baselineskip}
\subsection{3.ª resolución: Yang–Mills y la brecha del sector de curvatura}
El problema de Yang–Mills exige dos cosas inseparables: construir una teoría cuántica no trivial sobre \(\mathbb R^4\) para todo grupo compacto, conectado y simple del dominio considerado, y demostrar que su Hamiltoniano posee una brecha de masa positiva.\par
No basta con exhibir un operador finito con espectro positivo. Tampoco basta con escribir una acción de Wilson, una medida producto o un observable con un autovalor no nulo. La brecha debe pertenecer a la misma teoría gauge que posee localidad, positividad por reflexión, covariancia euclídea, reconstrucción de Wightman y curvatura física.\par
La característica gauge del continuo conserva ya la red refinable, la medida, las cuatro reflexiones, la holonomía, el color, la incidencia, el carry de fusión, la ruta, la memoria, el terminal y el lector de curvatura. El grupo \(G\) y el acoplamiento \(g\) entran después, mediante el entrelazador específico. No seleccionan retrospectivamente la genealogía.\par
\Needspace{6\baselineskip}
\subsubsection{Inequivalencia entre la fibra simple y la dinámica producto}
La primera distinción defensiva separa dos generadores que una notación antigua podía confundir.\par
El semigrupo de una sola fibra posee únicamente los niveles \(0\) y \(3\kappa_{\rm av}\). El generador producto completo integra, en cambio, todas las coordenadas independientes del estado:\par
\[
D_m^{\rm prod}
=
3\kappa_{\rm av}
\sum_{\iota\in\mathcal I_m}(I-E_{m,\iota}),
\]
donde las expectativas \(E_{m,\iota}\) son ortogonales y conmutantes.\par
Su espectro contiene\par
\[
0,\quad
3\kappa_{\rm av},\quad
6\kappa_{\rm av},\quad
9\kappa_{\rm av},\ldots
\]
según el número de coordenadas no constantes. Sustituir este generador por el de una fibra simple borraría las multiplicidades de Hoeffding y sería inmediatamente falsable: un vector dependiente de dos factores debe tener energía \(6\kappa_{\rm av}\), no \(3\kappa_{\rm av}\).\par
La brecha se demostrará para el generador producto y después se transportará al sector físico. No se obtiene confundiendo dos dinámicas diferentes.\par
\Needspace{6\baselineskip}
\subsubsection{Ausencia de coordenadas ocultas y exclusión de modos nulos espurios}
El estado completo se factoriza mediante un índice exhaustivo. En él aparecen una sola vez:\par
\begin{itemize}
\item marcos;
\item enlaces;
\item incidencias;
\item marcas;
\item puentes de refinamiento;
\item innovaciones nuevas de cada nivel.
\end{itemize}
Las subcolas infinitas se registran mediante sus cilindros escalares. Ninguna coordenada genealógica queda fuera de la familia de expectativas.\par
La descomposición de Hoeffding escribe el espacio como suma ortogonal de sectores \(\mathcal H_{m,S}\). Sobre cada uno,\par
\[
D_m^{\rm prod}u_S
=
3\kappa_{\rm av}|S|u_S.
\]
El sector \(S=\varnothing\) es la intersección de los rangos de todas las expectativas y está formado exactamente por las constantes:\par
\[
\bigcap_{\iota\in\mathcal I_m}
\operatorname{Ran}E_{m,\iota}
=
\mathbb C\Omega_m.
\]
Por tanto, el vacío es simple y todo vector ortogonal al vacío tiene energía al menos \(3\kappa_{\rm av}\).\par
Esto excluye la objeción de que la brecha provenga de no haber enumerado una coordenada. Si faltara un factor, existirían funciones no constantes invisibles para todas las expectativas y aparecerían cero-modos adicionales. La exhaustividad del índice impide exactamente ese fallo.\par
\Needspace{6\baselineskip}
\subsubsection{Realización observable de la brecha espectral}
Una cota inferior no basta para fijar el valor exacto de la brecha. Se necesita un testigo que alcance el umbral.\par
En cada collar de incidencia, la variable\par
\[
W_c(q)
=
\mathbf1_{\{q_1+\cdots+q_6=0\}}
-\frac13
\]
tiene media cero y satisface\par
\[
D_m^{\rm prod}W_c
=
3\kappa_{\rm av}W_c.
\]
Pero todavía debe demostrarse que \(W_c\) pertenece al sector gauge físico.\par
La acción de incidencia se escribe mediante la matriz\par
\[
(B^+x)_{ij}=x_i+x_j.
\]
Cada componente de \(x\in\mathbb F_3^4\) aparece tres veces al sumar las seis incidencias. En \(\mathbb F_3\),\par
\[
\sum_{i<j}(B^+x)_{ij}
=
3\sum_i x_i
=
0.
\]
Por tanto, la condición que define \(W_c\) es gauge invariante. El promedio sobre el grupo gauge completo conserva \(W_c\); no lo expulsa del espacio físico. Así, \(W_c\) es un autovector físico y el valor \(3\kappa_{\rm av}\) no es sólo una cota: es el primer nivel espectral alcanzado.\par
\Needspace{6\baselineskip}
\subsubsection{Construcción de la medida previa a la estimación de la brecha}
Una arista gruesa se refina en nueve incrementos ordenados. La medida fina no se inventa copiando nueve veces la medida gruesa. Se construye mediante el puente de calor condicionado por el producto de los nueve incrementos. Los tiempos de los subtramos pueden ser distintos y satisfacen únicamente que su suma sea el tiempo grueso.\par
La convolución del kernel de calor garantiza que el puente es una probabilidad. Una carta triangular separa entonces la variable gruesa de las innovaciones nuevas. La proyección hacia el nivel anterior multiplica los nueve incrementos y olvida exclusivamente las innovaciones.\par
De esta manera,\par
\[
(\widehat\pi_{m+1,m})_*
\widehat\mu_{m+1}
=
\widehat\mu_m.
\]
Las aristas compartidas por varias caras continúan siendo una misma variable. No se sustituyen por copias independientes. Carry y memoria impiden volver a contar en el nivel fino la cara que ya estaba presente en el nivel grueso.\par
Sobre esa familia proyectiva se construye primero el lector exacto de \(F^2\), su raíz positiva, la acción local y la densidad de Gibbs dependiente del acoplamiento:\par
\[
d\widehat\mu_{m,g}^{\rm YM}
=
Z_{m,g}^{-1}
e^{-S_{m,g}}\,
d\widehat\mu_m^0.
\]
Las ecuaciones de Schwinger–Dyson y Ward se calculan para esa medida. La brecha no se usa para elegir \(g\), para definir \(S_{m,g}\) ni para reponderar la ley. Aparece después, al estudiar el generador reversible asociado a la misma medida.\par
Ésta es la defensa contra la circularidad acción–ley–brecha: primero se construye la acción y su medida; después se demuestra su espectro.\par
\Needspace{6\baselineskip}
\subsubsection{Descenso del generador al cociente de calibre}
El espacio completo incluye enlaces, marcos, incidencias, color, orientación, ruta, memoria y terminal. La marginal que conserva únicamente los enlaces es una ablación posterior. No puede utilizarse para negar un observable que vive en la fibra de incidencia.\par
La proyección física se obtiene promediando la acción del grupo gauge completo. Como el generador integra coordenadas de forma equivariante,\par
\[
[\mathbb P_m^{\rm phys},D_m^{\rm prod}]=0.
\]
El semigrupo desciende, por tanto, al cociente físico. Su kernel se define primero sobre el estado completo y sólo después sobre las clases gauge. No se toma un operador restringido a una fibra y se lo declara kernel sobre todo el espacio.\par
El vacío simple y el autovector \(W_c\) sobreviven a la compresión. La brecha finita pertenece ya al espacio gauge-invariante.\par
\Needspace{6\baselineskip}
\subsubsection{Unicidad del proceso reconstruido a partir de \(4\) positividades por reflexión}
La teoría euclídea debe satisfacer positividad por reflexión en las cuatro direcciones. No se construyen cuatro medidas independientes. Una única medida de caras, obtenida del mismo kernel reversible, factoriza cada forma de Osterwalder–Schrader:\par
\[
\int
\overline{F(\Theta_{\nu}y)}F(y)\,d\Omega
=
\|\mathcal B_\nu F\|^2
\ge0,
\qquad \nu=1,\ldots,4.
\]
La positividad no es una hipótesis añadida. Es una norma cuadrada producida por Markov y reversibilidad.\par
Las marginales de caras opuestas recuperan el mismo semigrupo. Por tanto, las cuatro reconstrucciones OS poseen un único origen dinámico. Los conectores entre semirredes demostrarán después que reconstruyen el mismo espacio, el mismo vacío y el mismo Hamiltoniano.\par
\Needspace{6\baselineskip}
\subsubsection{Persistencia de la brecha espectral bajo el límite continuo}
Los mapas de refinamiento entrelazan las formas de Dirichlet, los semigrupos, el vacío y el testigo \(W_c\). El paso al continuo no se basa únicamente en observar que la brecha vale lo mismo en cada nivel. Esa igualdad, aislada, no impediría que apareciesen nuevos estados de energía arbitrariamente pequeña en el límite.\par
La construcción controla los productos cruzados mediante entrelazadores de bloque. Las coordenadas gruesas se distribuyen sobre sus descendientes con la normalización exacta. Las identidades de Gram demuestran que las publicaciones suavizadas forman familias de Cauchy. De este modo, la forma, el vacío y el autovector alcanzan el límite dentro del mismo sistema inductivo.\par
Se obtiene\par
\[
D_{\infty,g}^{\rm YM}
\succeq
3\kappa_{\rm av}(I-P_{\Omega_\infty}),
\]
y el testigo límite satisface\par
\[
D_{\infty,g}^{\rm YM}W_c
=
3\kappa_{\rm av}W_c.
\]
\Needspace{5\baselineskip}
Por tanto,\par
\[
\Delta_{\rm YM}^{\rm HMT}
=
3\kappa_{\rm av}>0.
\]
La desigualdad y el testigo convergen juntos. La brecha no se deduce de una mera semicontinuidad inferior ni se extrapola desde una malla aislada.\par
\Needspace{6\baselineskip}
\subsubsection{Pertenencia del testigo espectral al sector de curvatura}
Todavía quedaría una objeción si \(W_c\) fuese sólo una excitación combinatoria de un factor interno, desacoplada de la curvatura Yang–Mills.\par
La holonomía de incidencia elimina esa posibilidad. La suma de las seis incidencias determina una holonomía local de orden tres; \(W_c\) es una función de clase de esa holonomía. Por tanto pertenece al álgebra local gauge-invariante generada por los lazos.\par
El lector \(F^2\) se polariza para producir un espacio de curvaturas locales. Los entrelazadores de refinamiento transportan esas curvaturas conservando el producto de \(\mathfrak g\). Las sumas celulares convergen a una distribución gauge-covariante\par
\[
\mathbf F_{\mu\nu}
\in
\mathcal S'(\mathbb R^4;\mathfrak g)
\widehat\otimes
L^2(\widehat X_\infty,\widehat\mu_{\infty,g}^{\rm YM}).
\]
El sector cíclico generado por \(\mathbf F\) es reductor para el Hamiltoniano y contiene el testigo \(W_c\). En consecuencia, la brecha exacta pertenece al sector físico de curvatura. No es el espectro de un ruido auxiliar.\par
\Needspace{6\baselineskip}
\subsubsection{Localidad euclídea \(E(4)\) y reconstrucción continua de Osterwalder–Schrader}
El generador se descompone en circuitos locales de bloques con soporte uniformemente acotado. Una coloración nonádica separa los collares que pueden actualizarse simultáneamente. El producto de los circuitos preserva la identidad de Bianchi en cada paso, no sólo al finalizar una barrida.\par
Las traslaciones y rotaciones euclídeas actúan por relabelado de la red, los marcos, las incidencias y las holonomías. La continuidad fuerte no se prueba únicamente para un primer caos abstracto. Se controla primero sobre observables locales, luego sobre holonomías, plaquetas y clover; el suavizado de órbitas extiende la acción a un núcleo denso.\par
El sistema cofinal produce una red local \(\mathfrak A_{\rm YM}(O)\), un vacío \(\Omega_{\infty,g}\), una representación fuerte de \(E(4)\) y funcionales de Schwinger. Las pilas Markov y las transferencias OS se identifican con el mismo semigrupo continuo. La positividad por reflexión, la regularidad, la covariancia y el cluster exponencial permiten la reconstrucción de Osterwalder–Schrader.\par
La continuación produce campos locales de Wightman y una representación de Poincaré con espectro en el cono futuro. Las cuatro reconstrucciones euclídeas coinducen el mismo Hamiltoniano de curvatura.\par
La brecha controla además el cluster:\par
\[
|S^{\rm YM}(A\,\tau_rB)|
\le
e^{-3\kappa_{\rm av}r}
\|A\Omega_\infty\|
\|B\Omega_\infty\|.
\]
Así, el valor espectral positivo posee una consecuencia física continua y local; no queda encerrado en el cálculo finito.\par
\Needspace{6\baselineskip}
\subsubsection{Reconstrucción de la acción clásica como salida del sistema euclídeo}
La densidad finita gobernante es el lector exacto de \(F^2\). El clover no la sustituye. Funciona como carta suave para reconocer su límite.\par
Para una conexión suave, la expansión de Magnus de la holonomía de una plaqueta empieza por\par
\[
\log\operatorname{Hol}_{p_{\mu\nu}}
=
a_m^2F_{\mu\nu}
+\text{término impar}
+O(a_m^4).
\]
El promedio sobre las cuatro orientaciones del clover cancela el término impar. La suma de Riemann da\par
\[
S_{m,g}(A)
\longrightarrow
\frac1{4g^2}
\int_{\mathbb R^4}\|F_A(x)\|^2\,d^4x.
\]
Por tanto, la ley construida es reconocida como Yang–Mills porque su acción finita, sus holonomías y su curvatura convergen al funcional clásico correcto. No se define “teoría Yang–Mills” como cualquier teoría que posea el gap deseado.\par
La misma medida satisface Schwinger–Dyson y Ward. La relación entre pequeños lazos y curvatura, la identidad de Bianchi y la covariancia gauge se conservan en el límite. Sólo después de toda esta identificación se aplica la equivalencia espectral que transporta el gap.\par
\Needspace{6\baselineskip}
\subsubsection{Interpretación dimensional del parámetro de masa obtenido por el entrelazador terminal}
La brecha energética se publica dimensionalmente como\par
\[
m_{\rm gap}^{\rm HMT}
=
\frac{\hbar_{\rm int}}{c_{\rm int}}
\,3\kappa_{\rm av}>0.
\]
Este resultado no asigna una masa de Proca a un gluón elemental. La invariancia gauge permanece intacta. La masa corresponde al coste mínimo de excitar el sector colectivo gauge-invariante de curvatura por encima del vacío.\par
Esta distinción no rebaja el resultado: es exactamente la formulación requerida por el problema de Yang–Mills. La brecha pertenece al espectro físico de estados gauge-invariantes, no a un término de masa elemental añadido al lagrangiano.\par
\Needspace{6\baselineskip}
\subsubsection{Existencia de una brecha espectral positiva y reconstrucción masiva del sector de Yang–Mills}
Para todo grupo de Lie compacto, conectado y simple, la característica gauge
del continuo, seguida de su entrelazador específico, construye una teoría
cuántica no trivial de Yang--Mills sobre \(\mathbb R^4\).\par
La teoría posee:\par
\begin{itemize}
\item una familia proyectiva de medidas locales dependientes del acoplamiento;
\item una acción exacta reconocida en el límite como \(\frac1{4g^2}\int\|F\|^2\);
\item un cociente gauge físico;
\item positividad por reflexión en cuatro direcciones;
\item una red local y una acción fuerte de \(E(4)\);
\item funciones de Schwinger y continuación de Wightman;
\item vacío simple;
\item ecuaciones de Schwinger–Dyson, Ward y Bianchi;
\item un sector cíclico de curvatura;
\item la brecha exacta
  \[
  \Delta_{\rm YM}^{\rm HMT}=3\kappa_{\rm av}>0.
  \]
\end{itemize}
La conclusión resiste las objeciones previsibles porque no identifica el generador de una fibra con el producto; no deja coordenadas fuera del índice; no confunde la marginal de enlaces con el cociente físico; no extrapola la brecha desde una malla; no coloca el testigo fuera del sector de curvatura; no presupone continuidad euclídea; y no usa el gap para construir la acción o la medida.\par
Navier–Stokes demuestra que el continuo puede transportar una no linealidad sin perder la memoria necesaria para impedir una singularidad. Yang–Mills demuestra que el mismo continuo puede transportar una libertad gauge sin perder la curvatura necesaria para separar el vacío de la primera excitación física.\par

\Needspace{6\baselineskip}
\section{Clausuras algebraicas del continuo y límites de la reconstrucción extensional}
Las cinco construcciones \texttt{sp}, \texttt{sol}, \texttt{gau}, \texttt{coh} y \texttt{det} no son ramas que se añadan al continuo para resolver cinco problemas conocidos. Son características simultáneas del mismo estado enriquecido. Hodge y Birch–Swinnerton–Dyer son las dos publicaciones finales de esta arquitectura: una convierte una historia cohomológica compatible en una clase algebraica; la otra demuestra que las publicaciones analítica y aritmética de una misma línea determinantal coinciden.\par
\Needspace{6\baselineskip}
\subsection{4.ª resolución: Hodge y la construcción efectiva de la preimagen algebraica}
El problema de Hodge pregunta si toda clase racional de tipo \((p,p)\),\par
\[
\alpha\in
H^{2p}(X,\mathbb Q)\cap H^{p,p}(X),
\]
para una variedad proyectiva lisa compleja \(X\), procede de una clase algebraica\par
\[
\beta_\alpha\in \operatorname{CH}^{p}(X)_{\mathbb Q}
\]
tal que\par
\[
\operatorname{cl}_{X,p}(\beta_\alpha)=\alpha.
\]
La dificultad no consiste en escribir esta última igualdad. Consiste en construir \(\beta_\alpha\) sin haber introducido antes, bajo otro nombre, la existencia de la preimagen que se quiere demostrar.\par
La característica cohomológica \texttt{coh} proporciona el aparato anterior a la entrada de \(X\), \(p\) y \(\alpha\): historias supervivientes, refinamientos, relaciones, orientación, carry, memoria, frontera, ruta, multisección, terminal y lector. El dato clásico no crea ese aparato ni selecciona retrospectivamente un subdominio. Sólo pregunta qué publica ese aparato cuando recibe las observaciones de \(\alpha\).\par
En cada profundidad \(N\) aparecen dos módulos:\par
\[
\mathscr S_N(X,p),
\]
que contiene los estados cohomológicos enriquecidos, y\par
\[
\mathscr H_N^p(X),
\]
que contiene las observaciones cohomológicas visibles en esa profundidad. El lector finito\par
\[
e_N:\mathscr S_N(X,p)\longrightarrow\mathscr H_N^p(X)
\]
está coordinado con los mapas de restricción:\par
\[
\rho_{N+1,N}:\mathscr S_{N+1}\longrightarrow\mathscr S_N,
\qquad
q_{N+1,N}:\mathscr H_{N+1}^p\longrightarrow\mathscr H_N^p,
\]
mediante el cuadrado\par
\[
e_N\rho_{N+1,N}
=
q_{N+1,N}e_{N+1}.
\]
Esto impide que una observación admitida en una profundidad sea incompatible con su sombra anterior.\par
Para la clase \(\alpha\), la fibra finita que debe poblarse es\par
\[
\mathscr F_N(\alpha)
=
\{s\in\mathscr S_N(X,p):e_N(s)=q_N\alpha\}.
\]
Definirla no demuestra que sea no vacía. Tampoco demuestra que un elemento de \(\mathscr F_N(\alpha)\) pueda prolongarse hasta \(\mathscr F_{N+1}(\alpha)\). La resolución consiste precisamente en construir ambas cosas.\par
\Needspace{6\baselineskip}
\subsubsection{Construcción de la corrección relativa previa a la clase de cohomología}
Al refinar de \(N\) a \(N+1\) aparecen dos núcleos relativos:\par
\[
\mathscr K_{N+1}^{\mathrm S}
=
\ker\rho_{N+1,N},
\qquad
\mathscr K_{N+1}^{\mathrm H}
=
\ker q_{N+1,N}.
\]
El primero contiene las correcciones nuevas del estado que desaparecen al volver al nivel anterior. El segundo contiene las observaciones nuevas que tampoco eran visibles en el nivel anterior.\par
Las cartas nuevas, los cruces de Mayer–Vietoris, las orientaciones, el carry, la memoria, la frontera, la ruta y las coordenadas terminales producen una presentación finita. Sobre sus generadores se construyen dos mapas:\par
\[
a_{N+1}:
\mathscr A_{N+1}^{\mathrm{rel}}
\longrightarrow
\mathscr K_{N+1}^{\mathrm S},
\]
\[
b_{N+1}:
\mathscr A_{N+1}^{\mathrm{rel}}
\longrightarrow
\mathscr K_{N+1}^{\mathrm H},
\]
satisfaciendo\par
\[
e_{N+1}^{\mathrm{rel}}a_{N+1}=b_{N+1}.
\]
El mapa \(a_{N+1}\) produce una corrección del estado; \(b_{N+1}\) publica exactamente su corrección cohomológica.\par
Las observaciones relativas se ordenan por profundidad, soporte, hoja, ruta, carry, memoria y terminal. Con ese orden, la matriz de \(b_{N+1}\) es unitriangular:\par
\[
b_{N+1}(c_\mu)
=
\bar c_\mu+
\sum_{\nu<\mu}
B_{\nu\mu}\bar c_\nu.
\]
La diagonal unitaria no se elige para representar una clase concreta. Procede de la incidencia principal de cada generador; los términos inferiores quedan impuestos por las relaciones, la orientación y el carry. La tabla se construye sobre todas las observaciones relativas antes de introducir \(\alpha\).\par
La eliminación triangular produce explícitamente un inverso derecho:\par
\[
\sigma_{N+1}^{\mathrm{rel}}(\bar c_{\mu_0})=c_{\mu_0},
\]
\[
\sigma_{N+1}^{\mathrm{rel}}(\bar c_\mu)
=
c_\mu-
\sum_{\nu<\mu}
B_{\nu\mu}
\sigma_{N+1}^{\mathrm{rel}}(\bar c_\nu),
\]
y, por inducción,\par
\[
b_{N+1}\sigma_{N+1}^{\mathrm{rel}}=I.
\]
\Needspace{5\baselineskip}
Por tanto,\par
\[
s_{N+1}^{\mathrm{rel}}
=
a_{N+1}\sigma_{N+1}^{\mathrm{rel}}
\]
satisface\par
\[
e_{N+1}^{\mathrm{rel}}
s_{N+1}^{\mathrm{rel}}
=I.
\]
Éste es el cortafuegos contra la circularidad. La sección relativa no se supone, no se deduce de que la clase deseada ya sea algebraica y no se construye específicamente para \(\alpha\). Existe universalmente para toda observación relativa.\par
Tampoco se necesita que la presentación sea inyectiva. El núcleo puede contener expresiones celulares redundantes. La prueba requiere sobreyectividad hacia las observaciones relativas, no unicidad de la expresión. La desaparición del cokernel es consecuencia del inverso derecho; no es una condición gobernante introducida antes de la prueba.\par
\Needspace{6\baselineskip}
\subsubsection{El operador universal de extensión}
La estructura posee una degeneración unitaria\par
\[
\iota_{N+1,N}:\mathscr S_N\longrightarrow\mathscr S_{N+1},
\qquad
\rho_{N+1,N}\iota_{N+1,N}=I.
\]
Sea \(s\in\mathscr S_N\) y sea \(h\in\mathscr H_{N+1}^p\) una observación compatible:\par
\[
e_N(s)=q_{N+1,N}(h).
\]
El defecto de la prolongación trivial es\par
\[
\delta_{N+1}(s,h)
=
h-e_{N+1}\iota_{N+1,N}(s).
\]
La compatibilidad garantiza que\par
\[
\delta_{N+1}(s,h)
\in
\mathscr K_{N+1}^{\mathrm H}.
\]
Se define entonces\par
\[
\operatorname{Ext}_{N+1,N}(s,h)
=
\iota_{N+1,N}(s)
+
s_{N+1}^{\mathrm{rel}}
\bigl(\delta_{N+1}(s,h)\bigr).
\]
Por construcción:\par
\[
\rho_{N+1,N}
\operatorname{Ext}_{N+1,N}(s,h)
=s,
\]
\[
e_{N+1}
\operatorname{Ext}_{N+1,N}(s,h)
=h.
\]
La primera identidad conserva la historia anterior. La segunda realiza exactamente la nueva observación. La naturalidad de la tabla relativa garantiza que orientación, carry, memoria, ruta, frontera y terminal no se pierdan al refinar.\par
El operador funciona para todo par compatible \((s,h)\). No contiene una clase algebraica ni un ciclo elegido. Ésa es la razón precisa por la que su aplicación posterior a una clase de Hodge no presupone Hodge.\par
\Needspace{6\baselineskip}
\subsubsection{Construcción nivel a nivel de la historia global en el sistema inverso}
En profundidad inicial existe una sección basada\par
\[
s_0^{\mathrm{base}}:
\mathscr H_0^p\longrightarrow\mathscr S_0.
\]
Sus imágenes son estados semilla, no ciclos algebraicos elegidos para representar \(\alpha\). Sólo después se fija\par
\[
s_0=s_0^{\mathrm{base}}(q_0\alpha)
\]
y se define recursivamente\par
\[
s_{N+1}
=
\operatorname{Ext}_{N+1,N}
\bigl(s_N,q_{N+1}\alpha\bigr).
\]
Las identidades del operador universal prueban por inducción:\par
\[
e_Ns_N=q_N\alpha,
\qquad
\rho_{N+1,N}s_{N+1}=s_N.
\]
Por tanto, cada fibra \(\mathscr F_N(\alpha)\) es no vacía y la familia\par
\[
\xi_\alpha=(s_N)_{N\ge0}
\]
es una historia compatible:\par
\[
\xi_\alpha
\in
\varprojlim_N\mathscr F_N(\alpha).
\]
No se invoca un límite inverso abstracto para ocultar la falta de secciones. La secuencia se exhibe recursivamente. Tampoco se utiliza la anulación de \(\varprojlim^1\) como sustituto de la construcción: la sobreyectividad y la condición de Mittag–Leffler aparecen después como consecuencias de un operador de extensión ya materializado.\par
\Needspace{6\baselineskip}
\subsubsection{Emergencia de la clase algebraica como límite de la construcción relativa}
Las observaciones finitas y la coordenada terminal separan las publicaciones del límite. Como \(\xi_\alpha\) reproduce \(q_N\alpha\) en toda profundidad,\par
\[
\operatorname{ev}_{\infty,X,p}^{\mathrm{Hdg}}
(\xi_\alpha)
=
\alpha.
\]
Sólo entonces actúa la publicación clásica:\par
\[
X_\chi(X,p)
\xrightarrow{R_{\mathrm{Hdg}}}
K_0\!\left(
\operatorname{Perf}^{\ge p}(X)
\right)_{\mathbb Q}
\xrightarrow{\operatorname{ch}^{\mathrm{loc}}_p}
\operatorname{CH}^{p}(X)_{\mathbb Q}
\xrightarrow{\operatorname{cl}_{X,p}}
\operatorname{Hdg}^{p}(X)_{\mathbb Q}.
\]
Se define\par
\[
\beta_\alpha
=
\operatorname{ch}^{\mathrm{loc}}_p
\bigl(R_{\mathrm{Hdg}}(\xi_\alpha)\bigr).
\]
\Needspace{9\baselineskip}
La identidad de acción del lector,\par
\[
\operatorname{cl}_{X,p}
\operatorname{ch}^{\mathrm{loc}}_p
R_{\mathrm{Hdg}}
=
\operatorname{ev}_{\infty,X,p}^{\mathrm{Hdg}},
\]
da finalmente\par
\[
\operatorname{cl}_{X,p}(\beta_\alpha)
=
\alpha.
\]
Esta identidad del lector no prueba por sí sola que exista \(\xi_\alpha\). La existencia fue construida antes mediante la tabla unitriangular, el inverso derecho, el operador universal y la recursión. Mantener separadas esas dos etapas es justamente lo que impide que la demostración sea circular.\par
Para toda clase racional de tipo \((p,p)\), la característica cohomológica
construye la sección compatible \(\xi_\alpha\), el ciclo
\(\beta_\alpha\in\operatorname{CH}^p(X)_{\mathbb Q}\) y la identidad
\(\operatorname{cl}_{X,p}(\beta_\alpha)=\alpha\), según
\cref{thm:pdf2-hodge-extension-universal,thm:pdf2-hodge-completitud-final}.
La exhaustividad celular y FIN1 forman parte de la construcción probada del
inverso derecho y de su prolongación coinductiva; no son hipótesis externas de
promoción.\par
\Needspace{6\baselineskip}
\subsection{5.ª resolución: BSD y la coincidencia de las publicaciones analítica y aritmética}
La resolución de Birch–Swinnerton–Dyer no comienza intentando adivinar el orden de anulación de \(L(E,s)\), ni reuniendo retrospectivamente rango, regulador, grupo de Tate–Shafarevich y factores de Tamagawa.\par
La característica determinantal \texttt{det} ya conserva historia, producto fibrado, unidad común, orientación, carry, memoria, supervivencia, terminal, retículo, dualidad, raíz y publicaciones compatibles. Sólo después entra una curva elíptica \(E/\mathbb Q\).\par
\Needspace{6\baselineskip}
\subsubsection{Complejo integral global y descomposición en fibras primarias}
La construcción se realiza sobre un complejo libre integral de historias. Para cada primo \(\ell\) y cada potencia \(\ell^n\), los coeficientes \(E[\ell^n]\), las transiciones, Alexander–Whitney, el pullback de Manin, la unidad y las publicaciones se obtienen por cambio de base del mismo complejo integral.\par
La fibra ternaria no gobierna a las restantes. No se inventa una falsa extensión de \(\mathbb Z_3\) a \(\mathbb Z_\ell\). Esta precaución es imprescindible para que la conclusión comprenda toda \(\Sha(E)\) y todos los factores de Tamagawa, no únicamente su componente \(3\)-primaria.\par
La unidad común produce las publicaciones de Kato y Poitou–Tate. En cada bloque finita–singular se construye directamente la homotopía de relleno\par
\[
h_*=-vf,
\]
cuya frontera es la diferencia entre ambas publicaciones.\par
Las historias se ordenan usando primo, nivel, orientación, carry, memoria, frontera, terminal y localización. Se obtiene así una descomposición basada exhaustiva en:\par
\begin{itemize}
\item unidad raíz;
\item bloques finita–singular;
\item retículo local–global.
\end{itemize}
No queda una clase residual sin tipar.\par
La homotopía \(H_{\mathrm{FS}}\) no entra como premisa. En cada bloque normal se define su acción y se comprueba\par
\[
\partial H_\lambda+H_\lambda\partial=I.
\]
La suma transportada por la descomposición basada satisface\par
\[
\partial H_{\mathrm{FS}}
+
H_{\mathrm{FS}}\partial
=
p_{\mathrm{FS}}.
\]
El miembro derecho es el proyector sobre el sector finita–singular. No borra el retículo que conserva \(\Sha\) y Tamagawa. La homotopía reproduce la homotopía de relleno previamente construida; no se postula para hacer desaparecer una obstrucción.\par
\Needspace{6\baselineskip}
\subsubsection{Finitud del grupo aritmético previa a la evaluación del término principal}
El retículo residual posee una matriz integral \(D_E\). Antes de tomar su cokernel se construye un inverso racional mediante la parte diagonal terminal y la serie geométrica finita de la parte estrictamente triangular. Se verifican las dos composiciones.\par
\Needspace{5\baselineskip}
Por tanto,\par
\[
D_E\otimes\mathbb Q
\]
es un isomorfismo y\par
\[
F_E=\operatorname{coker}D_E
\]
es finito. Esta finitud se obtiene antes de usar rango, \(\Sha\), regulador o término principal.\par
La forma de Smith conserva todos los divisores elementales y recompone exactamente sus componentes primarias. No se deduce la finitud de \(\Sha\) de la fórmula final: se deriva posteriormente de su inclusión en un grupo finito ya construido.\par
\Needspace{6\baselineskip}
\subsubsection{Bockstein, Kato–FERS y Poitou–Tate}
El diferencial determinantal universal \(S_E(T)\) genera las páginas Bockstein. Su forma de Smith \(T\)-ádica determina\par
\[
d
=
\operatorname{ord}_{T=0}
\det S_E(T)
=
\sum_q \dim V_q.
\]
No se sustituyen las páginas por ese entero. Se conservan cada cociente, cada diferencial reducido, cada emparejamiento y cada jet. El coeficiente regular es\par
\[
R_{\mathrm{Boc}}^{\mathrm{der}}(S_E)
=
\left[
T^{-d}\det S_E(T)
\right]_{T=0}
\neq0.
\]
La equivalencia Kato–FERS se construye fila por fila. Alexander–Whitney corta cada historia en todos sus puntos. El corte terminal conserva el generador con peso cero; los demás cortes adquieren pesos positivos determinados por profundidad y carry. La tabla exhaustiva tiene la forma\par
\[
I+TB_i(T),
\]
con inversa convergente y determinante alternado igual a uno en \(T=0\). De este modo, la unidad, las hojas, la memoria, la orientación y el terminal se preservan antes de consultar ningún valor de \(L(E,s)\).\par
Poitou–Tate enlaza después las páginas Bockstein con Mordell–Weil. Los mapas de Kummer, publicación y pegado se escriben en ambos sentidos y satisfacen unidad y counidad. Son inversos antes de comparar dimensiones.\par
El desdoblamiento de jets conserva las \(q\) direcciones correspondientes a una fila de profundidad \(q\). No se pretende obtener \(q\) covectores de una única coordenada ni se completan dimensiones con identidades artificiales. El Gram conserva todos los cruces entre filas, profundidades y jets.\par
Sólo después se toman dimensiones:\par
\[
\operatorname{rank}E(\mathbb Q)
=
\sum_q\dim V_q
=
d.
\]
El rango y el orden determinantal coinciden porque ambos son publicaciones de las páginas construidas, no porque uno se haya definido copiando al otro.\par
\Needspace{6\baselineskip}
\subsubsection{Obtención del grupo de Tate–Shafarevich y de los factores de Tamagawa mediante la sucesión exacta global}
Del cono local reducido se construye una inyección\par
\[
\Sha(E)\hookrightarrow F_E.
\]
Los levantamientos estructurales de las componentes locales producen la proyección y su sección. La comprobación a nivel de cociclos da:\par
\[
0
\longrightarrow
\Sha(E)
\longrightarrow
F_E
\longrightarrow
\bigoplus_v\Phi_v(\mathbb F_v)
\longrightarrow
0.
\]
Como \(F_E\) ya era finito,\par
\[
|\Sha(E)|<\infty
\]
y\par
\[
|F_E|
=
|\Sha(E)|
\prod_v c_v.
\]
La exactitud no se deduce por comparación de cardinales. Primero se construyen los mapas y se identifican sus núcleos e imágenes; sólo después se toman órdenes.\par
\Needspace{6\baselineskip}
\subsubsection{Construcción independiente del miembro analítico de BSD}
La modularidad de \(E\) proporciona la forma nueva asociada. Su transformada de Mellin se divide en dos intervalos, y la ecuación funcional transforma la segunda mitad en una integral convergente con todas sus derivadas. Esto produce un germen holomorfo de \(L(E,s)\) alrededor de \(s=1\) sin suponer su orden de anulación.\par
En el semiplano de convergencia absoluta, los factores de Euler se realizan por operadores normales. Los truncados convergen en norma traza y su determinante de Fredholm recupera el producto de Euler. La reducción de Schur sobre las coordenadas terminales y un atlas de cartas determinantales transportan esa familia hasta el diferencial \(S_E\).\par
El resultado es una factorización construida:\par
\[
L(E,s)
=
u_E(s)
\det S_E(T_E(s)).
\]
\Needspace{7\baselineskip}
La carta satisface\par
\[
T_E(1)=0,
\qquad
T'_E(1)=1,
\]
y la unidad de transición satisface\par
\[
u_E(1)=1.
\]
Esta normalización se demuestra mediante las transiciones del atlas, el complemento Fredholm y los cambios de base. No se ajusta usando el rango, el regulador, \(\Sha\) ni el coeficiente final.\par
En consecuencia,\par
\[
L(E,s)
=
(s-1)^dU_E^{\mathrm{lead}}(s),
\]
con\par
\[
U_E^{\mathrm{lead}}(1)
=
R_{\mathrm{Boc}}^{\mathrm{der}}(S_E)
\neq0.
\]
\Needspace{5\baselineskip}
Por tanto,\par
\[
\operatorname{ord}_{s=1}L(E,s)=d.
\]
Combinándolo con Poitou–Tate:\par
\[
\operatorname{ord}_{s=1}L(E,s)
=
\operatorname{rank}E(\mathbb Q).
\]
\Needspace{6\baselineskip}
\subsubsection{Derivación interna del regulador a partir del emparejamiento de alturas}
La base integral de Mordell–Weil procede del sistema completo de jets. Sobre ella se construye la altura de Néron–Tate y su descomposición local–global.\par
El Gram contiene todos los cruces:\par
\[
G_E^{\mathrm{jet}}
=
\bigl(
\langle P_{q,a},P_{q',a'}\rangle_{\mathrm{NT}}
\bigr).
\]
No se supone diagonalidad entre profundidades ni se sustituyen direcciones adicionales por bloques identidad. Su determinante define\par
\[
\operatorname{Reg}(E)
=
\det G_E^{\mathrm{jet}}
>0.
\]
La extensión de escalares a \(\mathbb R\) se realiza antes de multiplicar por el regulador. Por ello el comparador entre las líneas Bockstein y Néron está bien tipado; no se intenta introducir un número real dentro de una línea todavía definida únicamente sobre \(\mathbb Q\).\par
\Needspace{6\baselineskip}
\subsubsection{Diagrama de \(4\) morfismos y evaluación del término principal}
La comparación final utiliza cuatro flechas de líneas determinantales:\par
\begin{enumerate}
\item Kato–FERS transporta la sección analítica.
\item Poitou–Tate la lleva hasta Mordell–Weil y su altura.
\item La sucesión exacta separa \(\Sha\) y Tamagawa.
\item La evaluación incorpora torsión, período de Néron y componentes locales.
\end{enumerate}
Cada flecha tiene dominio, codominio y acción sobre una base. Antes de introducir el elemento analítico se fijan dos trivializaciones independientes: la analítica y la de Néron.\par
La conmutatividad del cuadrado final procede de la conservación de las bases estructurales. No se impone igualando previamente los números que deben aparecer.\par
Aplicando la composición al elemento analítico se obtiene:\par
\[
\operatorname{ord}_{s=1}L(E,s)
=
\operatorname{rank}E(\mathbb Q)
=
r,
\]
y\par
\[
\frac{L^{(r)}(E,1)}
{r!\,\Omega_E}
=
\frac{
|\Sha(E)|
\operatorname{Reg}(E)
\prod_v c_v
}{
|E(\mathbb Q)_{\mathrm{tors}}|^2
}.
\]
Ésta es la fórmula completa de Birch--Swinnerton--Dyer para toda curva
elíptica \(E/\mathbb Q\), como establece
\cref{thm:pdf2-bsd-publicacion-final}. La familia compatible
\(\nu_{E,\ell,n}\) se construye dentro del complejo integral global y aporta
el comparador que enlaza sus fibras primarias; forma parte de la demostración
universal y no constituye una condición externa de promoción.\par
Su defensa contra la circularidad es estrictamente cronológica:\par
\begin{itemize}
\item el complejo integral se fija antes de separar primos;
\item la homotopía de relleno se construye antes de \(H_{\mathrm{FS}}\);
\item Smith se obtiene antes del rango;
\item Kato–FERS se verifica antes del determinante principal;
\item Poitou–Tate se construye antes de comparar dimensiones;
\item la exactitud \(\Sha\)–Tamagawa se prueba antes de tomar órdenes;
\item la continuación Mellin–Fredholm se construye antes del orden analítico;
\item la altura se construye antes del regulador;
\item las cuatro flechas se fijan antes de introducir el elemento analítico;
\item las trivializaciones se fijan antes de evaluar ambos lados.
\end{itemize}
Ningún elemento se normaliza usando el valor que después debe demostrar.\par
\Needspace{6\baselineskip}
\subsection{Factorización de los entrelazadores terminales RH–NS–YM–Hodge–BSD a través del continuo genealógico común}
Riemann, Navier–Stokes, Yang–Mills, Hodge y BSD no son cinco traducciones arbitrarias de un mismo vocabulario. Cada problema aparece cuando una publicación clásica olvida una parte diferente del estado enriquecido:\par
\begin{itemize}
\item Riemann olvida la coligación común que conserva simultáneamente los canales primo, gamma y polar.
\item Navier–Stokes olvida la columna causal completa de terminal, pérdidas, carry y memoria.
\item Yang–Mills olvida la coordinación entre incidencia, curvatura, ley, dinámica y sector físico.
\item Hodge conserva una clase cohomológica, pero no la historia compatible que produce su preimagen algebraica.
\item BSD separa la publicación analítica de la aritmética y pierde la línea determinantal común que las transporta.
\end{itemize}
La arquitectura HMT no fuerza cinco soluciones independientes. Conserva antes aquello que las formulaciones clásicas publican por separado:\par
\[
\begin{aligned}
\text{estado enriquecido}
&\longrightarrow \text{refinamientos compatibles}
\longrightarrow \text{supervivencia y terminal}\\
&\longrightarrow \text{continuo discreto}
\longrightarrow \text{publicaciones clásicas}.
\end{aligned}
\]
La razón por la que las soluciones funcionan no es una semejanza verbal. Es que, en los cinco casos, el teorema decisivo procede de una identidad de conservación material:\par
\begin{itemize}
\item conservación de energía o defecto;
\item conservación de historia bajo refinamiento;
\item conservación de la acción y del producto superficial;
\item conservación de la observación durante la extensión;
\item conservación de la unidad y de la base en líneas determinantales.
\end{itemize}
Por eso los problemas clásicos deben aparecer editorialmente al final. Primero se construye el objeto común; después HMT–MD desarrolla sus consecuencias; luego Mecánica Dimensional realiza sus transductores; sólo finalmente se leen las cinco conclusiones en sus lenguajes clásicos.\par
\Needspace{6\baselineskip}
\subsection{Distinción tipada entre los cinco cierres terminales y el corredor lógico ZFC--ciclo}
La incidencia con Zermelo--Fraenkel, ZFC y la construcción de von Neumann debe formularse con precisión y en el orden causal del corpus. Las cinco resoluciones terminales y el corredor lógico ZFC--ciclo son demostraciones distintas, asentadas sobre la misma estructura discreta del continuo: el alcance local de una de ellas no puede operar como negación transversal de la otra. Los cinco cierres establecen sus respectivos resultados clásicos; el corredor lógico específico conserva, además, la cadena\par
\[
\begin{aligned}
 \mathcal C_{\mathrm{cont}}^{\mathrm{disc}}
 &\longrightarrow\mathcal L_0\xrightarrow{\Phi}\mathcal L_1
 \longrightarrow\gamma_F\longrightarrow w(\gamma_F)\\
 &\longrightarrow\mathsf{NoEscape}_{24}
 \longrightarrow\mathsf{AntiDiag}_{369}\\
 &\longrightarrow\mathsf{G\ddot{o}delIII}_{\mathrm{HMT}}
 \longrightarrow\mathsf{CycleCriterion}_{\mathrm{ZFC}},
\end{aligned}
\]
y publica el criterio fuerte recuperado\par
\[
 \boxed{
 \neg\operatorname{Consis}(\mathrm{ZFC})
 \Longleftrightarrow
 \exists C\;\bigl(C\text{ es un ciclo fuera de banda y }|C|>24\bigr).}
\]
Esta equivalencia posee su propio dominio, sus codificadores y su prueba; no se deduce de los cinco teoremas terminales ni queda limitada por una observación local acerca de cualquiera de ellos. Del mismo modo, la resolución cardinal ya construida conserva simultáneamente\par
\[
 \mathsf{CH}_{\mathrm{HMT}},
 \qquad
 M_h=L[h]\models 2^{\aleph_0}=\aleph_1,
\]
sin convertir la recta real o la sentencia extensional en entradas del generador.\par

Una representación extensional formulada en el lenguaje de ZFC puede codificar los portadores conjuntistas de los objetos empleados:\par
\begin{itemize}
\item un estado finito es un conjunto;
\item una historia es una sucesión;
\item un árbol es un conjunto de palabras;
\item una transición es una función o relación;
\item un sistema inverso es una familia conjuntista;
\item su límite es un subconjunto de un producto.
\end{itemize}
Esa codificación posterior no contiene por sí sola la identidad genealógica del objeto ni gobierna los teoremas construidos antes de la publicación extensional. La insuficiencia aparece cuando el reducto extensional se trata como si agotara ruta, orientación, acarreo, memoria, frontera y terminal.\par
El axioma de extensionalidad afirma:\par
\[
\forall A\,\forall B
\left[
\forall x\,
(x\in A\leftrightarrow x\in B)
\Rightarrow A=B
\right].
\]
La afirmación es correcta para conjuntos. Dos conjuntos con los mismos elementos son el mismo conjunto.\par
Pero de ella no se sigue que una publicación extensional permita reconstruir la historia que fue olvidada al producirla.\par
Sea\par
\[
\operatorname{Pub}:
\mathscr X_{\mathrm{gen}}
\longrightarrow
\mathscr X_{\mathrm{ext}}
\]
\Needspace{6\baselineskip}
un lector que envía un estado genealógico a su valor extensional. Puede ocurrir que\par
\[
\operatorname{Pub}(x)
=
\operatorname{Pub}(y)
\]
mientras\par
\[
x\neq y
\]
como historias, porque poseen distinta ruta, orientación, carry, memoria, frontera o terminal.\par
El reducto identifica las dos imágenes en \(\mathscr X_{\mathrm{ext}}\), pero esa igualdad de codominio no implica \(x=y\) en el dominio enriquecido. El error de tipo consiste en convertir la igualdad de publicaciones en igualdad genealógica.\par
Puede definirse la relación\par
\[
x\sim_{\mathrm{ext}}y
\quad\Longleftrightarrow\quad
\operatorname{Pub}(x)
=
\operatorname{Pub}(y).
\]
Entonces, en el dominio correspondiente,\par
\[
\mathscr X_{\mathrm{ext}}
\simeq
\mathscr X_{\mathrm{gen}}/
{\sim_{\mathrm{ext}}}.
\]
La fibra\par
\[
\operatorname{Pub}^{-1}(a)
\]
contiene las genealogías que producen la misma sombra \(a\). Extensionalidad describe correctamente el cociente; no proporciona una sección canónica que reconstruya una historia concreta de cada fibra.\par
Ni siquiera el axioma de elección resuelve esta carencia. Elección puede seleccionar un representante de cada fibra bajo sus hipótesis, pero:\par
\begin{itemize}
\item no lo hace canónico;
\item no conserva necesariamente operaciones;
\item no garantiza naturalidad bajo refinamiento;
\item no reconstruye carry, memoria o terminal;
\item no demuestra que la selección corresponda a la genealogía física o matemática correcta.
\end{itemize}
Elegir un representante no equivale a recuperar su procedencia.\par
\Needspace{6\baselineskip}
\subsubsection{El ejemplo de los ordinales de von Neumann}
La construcción de von Neumann\par
\[
0=\varnothing,
\qquad
n+1=n\cup\{n\}
\]
es exacta. El ordinal\par
\[
n=\{0,1,\ldots,n-1\}
\]
codifica perfectamente su orden extensional.\par
\Needspace{6\baselineskip}
Pero el mismo número puede aparecer como salida de historias distintas:\par
\[
6=3+3=2\cdot3=7-1.
\]
Como ordinal de von Neumann, el resultado es el mismo conjunto. Eso es correcto: la aritmética necesita que las cuatro expresiones publiquen el mismo \(6\).\par
Lo que el ordinal desnudo no conserva es qué operación, ruta, acarreo u orientación produjo esa aparición concreta. Si esa procedencia importa para una aplicación posterior, debe mantenerse como parte del objeto enriquecido:\par
\[
\widehat n_x
=
\{(k,x_k):k<n\}.
\]
La proyección\par
\[
(k,x_k)\longmapsto k
\]
recupera exactamente el ordinal de von Neumann. HMT no elimina el ordinal; lo reconoce como publicación de una estructura que puede contener más información.\par
La jerarquía acumulativa \(V_\alpha\) conserva pertenencia y rango. No registra automáticamente una causalidad interna de construcción diferente de la codificación elegida. Reintroducir etiquetas, árboles o categorías permite transportar datos genealógicos adicionales, pero ese enriquecimiento confirma precisamente que la genealogía no estaba contenida en el valor extensional después de haberla olvidado.\par
\Needspace{6\baselineskip}
\subsubsection{La recta real como publicación posterior}
La misma distinción afecta al continuo. Un número real puede construirse como corte de Dedekind, clase de sucesiones de Cauchy u otro representante extensional equivalente. Esa equivalencia es matemáticamente correcta.\par
HMT añade una pregunta anterior: ¿qué sistema compatible de estados, refinamientos y terminales publica ese real?\par
La recta\par
\[
\mathbb R
\]
aparece después mediante una evaluación arquimediana. Conserva el valor publicado, pero no necesariamente la historia completa que lo produjo. Dos estados pueden compartir la misma evaluación real sin compartir genealogía.\par
Por eso la estructura discreta del continuo no debe introducir \(\mathbb R\) como entrada. Primero se construye la historia compatible; después se publica su sombra arquimediana.\par
El entrelazador con la cardinalidad clásica se construye de manera explícita.
La codificación por bloques proporciona el homeomorfismo
\(\beta:\Omega_\Gamma^{\rm cal}\simeq(\mathbb F_3^6)^{\mathbb N}\); sus
aproximantes racionales construyen \(\mathbb R_{\rm HMT}\), y el cociente por
igualdad de valor publica \(\mathbb R\). El código genealógico \(h\) induce
el buen orden y determina el modelo canónico \(M_h=L[h]\), dentro del cual
la construcción prueba la sentencia clásica
\(2^{\aleph_0}=\aleph_1\). Los resultados de independencia de Gödel y Cohen
clasifican el lenguaje extensional de ZFC después de olvidar \(h\); no
refutan la decisión producida por la estructura enriquecida que conserva el
selector.\par
\Needspace{6\baselineskip}
\subsubsection{Distinción entre identidad genealógica, equivalencia observacional e igualdad extensional}
La carencia es tipológica: una codificación conjuntista añadida puede transportar una genealogía ya construida, mientras el reducto extensional desnudo carece del mapa que produjo la publicación. Capacidad de codificación y suficiencia fundacional son afirmaciones diferentes.\par
Deben distinguirse tres nociones:\par
\[
\text{identidad genealógica},
\qquad
\text{equivalencia observacional},
\qquad
\text{igualdad extensional}.
\]
La primera exige igualdad de historia, ruta, memoria, carry y terminal. La segunda exige igualdad respecto de un lector determinado. La tercera exige igualdad del objeto publicado en el codominio conjuntista.\par
Confundirlas genera falsos problemas de reconstrucción. Una sombra puede ser exacta y, al mismo tiempo, insuficiente para recuperar el objeto que la produjo.\par
Ésta es la aportación fundacional que las cinco resoluciones hacen visible. Cada problema clásico trabajaba con una publicación correcta, pero empobrecida:\par
\begin{itemize}
\item una forma explícita sin toda su genealogía común;
\item un campo sin la memoria causal completa;
\item una teoría gauge sin la coordinación íntegra de su proceso;
\item una clase cohomológica sin su historia de extensión;
\item una función \(L\) separada de su línea aritmética.
\end{itemize}
La estructura HMT mantiene esas correlaciones antes de publicar las sombras. Por eso no necesita reconstruir a posteriori, dentro de cada problema, información que nunca fue destruida.\par
La conclusión fundacional es la siguiente:\par
\begin{quote}
El cierre extensional es una publicación posterior del cierre genealógico y no puede sustituirlo. HMT construye la estructura enriquecida antes de la recta real, demuestra \(\mathsf{CH}_{\mathrm{HMT}}\), realiza la clausura ordinal en \(L[h]\) y establece separadamente el corredor Gödel--HMT III y el criterio ZFC--ciclo. Una codificación extensional posterior conserva únicamente aquello que el morfismo de publicación declara; no revoca ninguno de esos resultados.\par
\end{quote}
\Needspace{6\baselineskip}
\subsection{Factorización de RH–NS–YM–Hodge–BSD como publicaciones terminales del continuo genealógico común}
Los cinco problemas aparecen así como cinco lecturas terminales de un único borde. No se resuelven porque una analogía los reúna, sino porque el continuo enriquecido conserva desde el principio las correlaciones que cada formulación clásica había separado.\par
La dependencia causal definitiva se factoriza como:\par
\[
\begin{aligned}
\mathrm{APP}
&\longrightarrow \mathrm{TRIT}
\longrightarrow \mathrm{TPK}
\longrightarrow \widehat x\\
&\longrightarrow \text{cinco características consustanciales}
\longrightarrow \text{terminal y límite}
\longrightarrow \mathfrak C_{\mathrm{cont}}^{\mathrm{disc}}
\end{aligned}
\]
y sólo después:\par
\[
\mathfrak C_{\mathrm{cont}}^{\mathrm{disc}}
\longrightarrow
\begin{cases}
\text{positividad de Weil y RH},\\
\text{regularidad global suave de Navier--Stokes para todo dato periódico suave},\\
\text{Yang--Mills y brecha de masa para todo grupo compacto, conectado y simple},\\
\text{algebraicidad de toda clase racional de Hodge},\\
\text{igualdad completa de BSD para toda curva elíptica sobre }\mathbb Q.
\end{cases}
\]
Las cinco líneas finales no son cinco ramas internas del continuo. Son cinco publicaciones clásicas posteriores de cinco capacidades inseparables que ya coexistían en el mismo estado.\par
Ese es el cierre que la formulación anterior no podía todavía formular: no estamos ante cinco problemas desconectados que casualmente admiten un mismo vocabulario, sino ante cinco pérdidas de información distintas corregidas por una sola arquitectura que se negó a destruirla.\par
